\chapter{Mapa de Decisão --- Top 10 por Perfil de Relocação}\label{mapa-decisao}

Este capítulo traduz os dados GSS em recomendações práticas para cinco perfis distintos de relocação. Cada perfil pondera os critérios do GSS de forma diferente: quem prioriza segurança física valoriza mais o sub-critério B; quem prioriza autossuficiência, o sub-critério D; e assim por diante.

\medskip
As listas abaixo são resultado da ponderação dos sub-critérios GSS (A, B, C, D, E) com pesos ajustados por perfil, filtrando distritos com GSS$\geq$6,0 e dados completos. Todos os distritos listados possuem análise completa nas seções de departamento correspondentes.

\section*{Perfil 1 --- Família com Filhos em Idade Escolar}\label{perfil-familia}
\addcontentsline{toc}{section}{Perfil 1 --- Família com Filhos}

\textbf{Prioridades:} Segurança urbana (B), estabilidade institucional e serviços (E), baixo risco natural (C), infraestrutura de saúde e educação.

\textbf{Filtros aplicados:} B $\geq$ 6,0; E $\geq$ 6,5; distância de fronteiras problemáticas; presença de escola secundária e hospital regional.

\begin{center}
\small
{\renewcommand{\arraystretch}{1.3}
\begin{tabular}{clllc}
\toprule
\textbf{\#} & \textbf{Distrito} & \textbf{Departamento} & \textbf{Por que indica?} & \textbf{GSS} \\
\midrule
1 & Quiindy & Paraguarí & Cidade tranquila, hospital regional, GSS topo & 8,5 \\
2 & Acahay & Paraguarí & Interior seguro, serviços básicos, solo fértil & 8,1 \\
3 & Eusebio Ayala & Cordillera & Próximo à Asunção, ótima infraestrutura & 8,0 \\
4 & Colonia Independencia & Guairá & Comunidade organizada, clima ameno, escolas & 7,7 \\
5 & Hernandarias & Alto Paraná & Cidade consolidada, Itaipu próxima, serviços & 8,0 \\
6 & Caaguazú & Caaguazú & Capital regional, hospital, universidade & 7,7 \\
7 & San Patricio & Misiones & Interior tranquilo, baixa densidade, seguro & 7,7 \\
8 & San Pedro del Ycuamandyyú & San Pedro & Capital departamental, serviços completos & 7,7 \\
9 & Itacurubi del Rosario & San Pedro & Interior pacífico, custo de vida baixíssimo & 8,0 \\
10 & Alto Vera & Itapúa & Colônia organizada, infraestrutura sólida & 8,0 \\
\bottomrule
\end{tabular}}
\end{center}

\textit{Nota: evitar distritos próximos à fronteira leste (Alto Paraná/Amambay) e ao norte (Concepción) para este perfil.}

\clearpage

\section*{Perfil 2 --- Empreendedor e Investidor}\label{perfil-empreendedor}
\addcontentsline{toc}{section}{Perfil 2 --- Empreendedor e Investidor}

\textbf{Prioridades:} Estabilidade institucional (E), logística e acesso a mercados, crescimento econômico, segurança do patrimônio, facilidade de abertura de empresa.

\textbf{Filtros aplicados:} E $\geq$ 7,0; acesso a rota principal; crescimento econômico documentado; conectividade 4G disponível.

\begin{center}
\small
{\renewcommand{\arraystretch}{1.3}
\begin{tabular}{clllc}
\toprule
\textbf{\#} & \textbf{Distrito} & \textbf{Departamento} & \textbf{Por que indica?} & \textbf{GSS} \\
\midrule
1 & Hernandarias & Alto Paraná & Zona franca Itaipu, hub agroindustrial & 8,0 \\
2 & Villarrica & Guairá & Capital cultural, base logística sul & 6,5 \\
3 & Encarnación & Itapúa & 2ª cidade do país, turismo, Mercosul & 7,2 \\
4 & Luque & Central & Polo tecnológico, aeroporto, FIFA-city & 6,8 \\
5 & San Lorenzo & Central & Polo universitário, indústria consolidada & 6,5 \\
6 & Coronel Oviedo & Caaguazú & Hub rodoviário central, crescimento rápido & 6,3 \\
7 & Caaguazú & Caaguazú & Agronegócio, soja, frigorifico, crescimento & 7,7 \\
8 & Santa Rosa del Monday & Alto Paraná & Soja e pecuária de alta produtividade & 8,2 \\
9 & Salto del Guairá & Canindeyú & Comércio fronteiriço, potencial logístico & 6,8 \\
10 & Filadelfia & Boquerón & Hub menonita, pecuária, exportação & 7,2 \\
\bottomrule
\end{tabular}}
\end{center}

\section*{Perfil 3 --- Agricultor e Buscador de Autossuficiência}\label{perfil-agricultor}
\addcontentsline{toc}{section}{Perfil 3 --- Agricultor e Autossuficiência}

\textbf{Prioridades:} Solo e recursos hídricos (D), baixo risco natural (C), distância de alvos e risco social baixo (A+B), acesso ao Aquífero Guarani.

\textbf{Filtros aplicados:} D $\geq$ 7,5; cobertura do Aquífero Guarani ou rios perenes; fertilidade comprovada; GSS $\geq$ 7,5.

\begin{center}
\small
{\renewcommand{\arraystretch}{1.3}
\begin{tabular}{clllc}
\toprule
\textbf{\#} & \textbf{Distrito} & \textbf{Departamento} & \textbf{Por que indica?} & \textbf{GSS} \\
\midrule
1 & Capitán Miranda & Itapúa & Solo top, Guarani acessível, clima ideal & 8,9 \\
2 & General Artigas & Itapúa & Fronteira Argentina, solo elite, seguro & 8,8 \\
3 & Itapúa Poty & Itapúa & Interior fértil, baixíssima densidade & 8,1 \\
4 & San José de los Arroyos & Caaguazú & Solo basáltico, água abundante, produtivo & 8,6 \\
5 & Katuete & Canindeyú & Solos profundos, Guarani acessível & 8,1 \\
6 & Santa Rosa del Monday & Alto Paraná & Soja mecanizável, alta produtividade & 8,2 \\
7 & San Alberto & Alto Paraná & Interior verde, Solo Oxissolo, produtivo & 8,0 \\
8 & Itacurubi del Rosario & San Pedro & Preço baixo, fertilidade boa, paz & 8,0 \\
9 & Villa del Rosario & San Pedro & Solo fértil, preço acessível, isolado & 8,0 \\
10 & Alto Vera & Itapúa & Colônia produtiva, solo argiloso fértil & 8,0 \\
\bottomrule
\end{tabular}}
\end{center}

\clearpage

\section*{Perfil 4 --- Aposentado de Baixo Custo}\label{perfil-aposentado}
\addcontentsline{toc}{section}{Perfil 4 --- Aposentado de Baixo Custo}

\textbf{Prioridades:} Segurança (B), custo de vida baixo, clima agradável, acesso a serviços de saúde básicos, comunidade de estrangeiros estabelecida.

\textbf{Filtros aplicados:} B $\geq$ 6,5; custo de vida abaixo da média nacional; hospital ou USF regional a menos de 50 km; altitude entre 100--400 m (clima temperado).

\begin{center}
\small
{\renewcommand{\arraystretch}{1.3}
\begin{tabular}{clllc}
\toprule
\textbf{\#} & \textbf{Distrito} & \textbf{Departamento} & \textbf{Por que indica?} & \textbf{GSS} \\
\midrule
1 & Quiindy & Paraguarí & Tranquilidade, custo baixo, ar puro & 8,5 \\
2 & Acahay & Paraguarí & Clima ameno, rural seguro, fácil integração & 8,1 \\
3 & Colonia Independencia & Guairá & Comunidade europeia, vinhos, clima suave & 7,7 \\
4 & San Bernardino & Cordillera & Lago Ypacaraí, comunidade internacional & 6,7 \\
5 & Caazapá & Caazapá & Interior pacífico, cultura local rica & 6,5 \\
6 & San Juan Bautista & Misiones & Capital tranquila, clima agradável & 6,5 \\
7 & Encarnación & Itapúa & Costanera famosa, turismo, comunidade BR & 7,2 \\
8 & Roque González de Santa Cruz & Paraguarí & Interior sereno, custo baixíssimo & 7,9 \\
9 & Eusebio Ayala & Cordillera & Próximo a Asunção, tranquilo, serviços & 8,0 \\
10 & Benjamin Aceval & Presidente Hayes & Chaco acessível, risco baixo, espaço & 8,2 \\
\bottomrule
\end{tabular}}
\end{center}

\section*{Perfil 5 --- Nômade Digital e Trabalhador Remoto}\label{perfil-nomade}
\addcontentsline{toc}{section}{Perfil 5 --- Nômade Digital}

\textbf{Prioridades:} Conectividade (4G/fibra), estabilidade elétrica, segurança urbana, custo de vida, qualidade de vida urbana, comunidade internacional.

\textbf{Filtros aplicados:} Cobertura 4G $\geq$ 70\%; fibra ótica disponível no município; GSS $\geq$ 6,0; serviços urbanos (coworking, cafés, bancos).

\begin{center}
\small
{\renewcommand{\arraystretch}{1.3}
\begin{tabular}{clllc}
\toprule
\textbf{\#} & \textbf{Distrito} & \textbf{Departamento} & \textbf{Por que indica?} & \textbf{GSS} \\
\midrule
1 & Asunción & Distrito Capital & Capital, melhor conectividade, bairros inter. & 7,7 \\
2 & Luque & Central & FIFA-city, fibra, comunidade tech & 6,8 \\
3 & San Lorenzo & Central & Universidades, conexão, custo razoável & 6,5 \\
4 & Encarnación & Itapúa & Qualidade de vida alta, fibra, costanera & 7,2 \\
5 & Hernandarias & Alto Paraná & Hub Itaipu, fibra, crescimento rápido & 8,0 \\
6 & Lambaré & Central & Bairro nobre de Asunção, fibra, seguro & 6,8 \\
7 & Coronel Oviedo & Caaguazú & Hub central, conectividade crescente & 6,3 \\
8 & Villarrica & Guairá & Cidade universitária, boa conectividade & 6,5 \\
9 & Salto del Guairá & Canindeyú & Comércio dinâmico, fibra disponível & 6,8 \\
10 & Filadelfia & Boquerón & Comunidade organizada, conectividade razoável & 7,2 \\
\bottomrule
\end{tabular}}
\end{center}

\clearpage

\section*{Comparativo dos Top 5 por Perfil}\label{comparativo-perfis}
\addcontentsline{toc}{section}{Comparativo por Perfil}

A tabela abaixo consolida os distritos que aparecem em mais de um perfil, indicando sua versatilidade:

\begin{center}
\small
{\renewcommand{\arraystretch}{1.3}
\begin{tabular}{lcccccl}
\toprule
\textbf{Distrito / Departamento} & \textbf{P1} & \textbf{P2} & \textbf{P3} & \textbf{P4} & \textbf{P5} & \textbf{GSS} \\
& \scriptsize Família & \scriptsize Invest. & \scriptsize Agric. & \scriptsize Aposen. & \scriptsize Digital & \\
\midrule
Hernandarias / Alto Paraná & $\checkmark$ & $\checkmark$ & & & $\checkmark$ & 8,0 \\
Encarnación / Itapúa & & $\checkmark$ & & $\checkmark$ & $\checkmark$ & 7,2 \\
Quiindy / Paraguarí & $\checkmark$ & & & $\checkmark$ & & 8,5 \\
Eusebio Ayala / Cordillera & $\checkmark$ & & & $\checkmark$ & & 8,0 \\
Itacurubi del Rosario / San Pedro & $\checkmark$ & & $\checkmark$ & & & 8,0 \\
Santa Rosa del Monday / Alto Paraná & & $\checkmark$ & $\checkmark$ & & & 8,2 \\
Villarrica / Guairá & & $\checkmark$ & & & $\checkmark$ & 6,5 \\
Caaguazú / Caaguazú & $\checkmark$ & $\checkmark$ & & & & 7,7 \\
Filadelfia / Boquerón & & $\checkmark$ & & & $\checkmark$ & 7,2 \\
Colonia Independencia / Guairá & $\checkmark$ & & & $\checkmark$ & & 7,7 \\
\bottomrule
\end{tabular}}
\end{center}

\medskip

\textbf{Distrito mais versátil:} \textbf{Hernandarias} (Alto Paraná) aparece em 3 dos 5 perfis, com GSS 8,0 e presença confirmada de fibra ótica, solo produtivo, hospital regional e polo industrial Itaipu. É a localidade com melhor equilíbrio geral entre os distritos analisados.

\textbf{Segundo lugar:} \textbf{Encarnación} (Itapúa) combina qualidade urbana, turismo, conectividade e proximidade à Argentina, sendo especialmente indicada para aposentados e nômades digitais que valorizam vida urbana agradável.

\medskip
Para análise detalhada de cada distrito, consulte a ficha completa na respectiva seção de departamento.

\clearpage
